\chapter{Firmware and Control}
\label{ch:firmware}

The firmware is written in C++ for the Arduino core for the ESP32~\cite{arduino_esp32} and uses
the MPU6050\_light library~\cite{mpu6050light} to read the IMU. This chapter explains its structure
and the control laws. The firmware is located in the folder \cmd{firmware/cubilox\_mk1/}.

\section{Source files}

\begin{table}[htbp]
  \centering
  \caption{Source files of the firmware.}
  \label{tab:files}
  \begin{tabularx}{\textwidth}{lX}
    \toprule
    File & Content \\
    \midrule
    \cmd{cubilox\_mk1.ino} & \cmd{setup()} and \cmd{loop()}: start-up, sensor filters, state
      machine, edge and vertex controllers \\
    \cmd{config.h} & pin definitions, controller gains, limits and timing constants, shared state
      variables \\
    \cmd{calibration.ino} & gyroscope calibration at start-up, edge and vertex calibration,
      storage in flash, pose detection \\
    \cmd{commands.ino} & command interface over USB and Bluetooth, storage of the gains \\
    \cmd{hardware.ino} & log output, battery measurement, buzzer, motor control, encoder
      interrupts, motor test \\
    \bottomrule
  \end{tabularx}
\end{table}

The Arduino IDE combines all \cmd{.ino} files of the sketch folder into one program; global
variables shared between the files are therefore declared in \cmd{config.h}.

\section{Start-up}

\cmd{setup()} performs the following steps:
\begin{enumerate}
  \item start the USB serial port (\SI{115200}{baud}) and Bluetooth with the device name
        \cmd{Cubilox\_Mk1},
  \item load the gains and the calibration from flash,
  \item configure the pins, attach the PWM outputs (\SI{20}{kHz}, 8~bit) and the encoder
        interrupts, stop all motors and engage the brake,
  \item initialise the IMU (gyroscope range \SI{\pm 500}{\degree/s}, accelerometer range
        \SI{\pm 2}{g}, the defaults of the library),
  \item calibrate the gyroscope (the cube must lie still, see \cref{sec:calibration}),
  \item measure the battery voltage and report the calibration status.
\end{enumerate}
After that, the cube waits until it is placed on a calibrated edge or vertex and starts balancing
automatically.

\section{Main loop and timing}

\cmd{loop()} runs continuously. In every pass it processes incoming commands, reads the IMU and
updates the sensor filters. The controllers themselves run whenever at least \SI{5}{ms} have
passed since the last control step, i.e.\ at up to \SI{200}{Hz}. The debug output reports the
longest gap between two control steps (\cmd{loop=}); on the real cube it is \SIrange{5}{6}{ms}.

The wheel speeds are computed from the encoder counts in a window of \SI{20}{ms} and scaled by 5,
so that their unit is \enquote{encoder counts per \SI{100}{ms}}. All wheel speeds and wheel
momenta in this chapter use this unit.

\section{State machine}
\label{sec:state_machine}

The firmware has three states:
\begin{description}
  \item[WAITING] The motors are off and braked. The firmware compares the measured acceleration
    vector with the calibrated edge and vertex poses. If it is closer than \SI{0.25}{g} to one of
    them, the cube is close enough to the balance point (edge: $|\theta| < \ang{1}$, vertex:
    tilt $< \ang{1.5}$), it is at rest ($\left||\vect{a}| - \SI{1}{g}\right| < \SI{0.3}{g}$), the
    battery is not empty and these conditions hold for \SI{400}{ms}, the brake is released and the
    corresponding balancing state is entered.
  \item[RUNNING\_EDGE] The edge controller drives motor X. The motor command is ramped up over
    \SI{800}{ms}. Balancing stops if the angle exceeds \ang{15} or the cube is moved
    ($\left||\vect{a}| - \SI{1}{g}\right| > \SI{0.3}{g}$).
  \item[RUNNING\_VERTEX] The vertex controller drives all three motors. The command is ramped up
    over \SI{150}{ms}. Balancing stops if the tilt exceeds \ang{25}, if a shock
    ($\left||\vect{a}| - \SI{1}{g}\right| > \SI{1.2}{g}$) lasts longer than \SI{150}{ms}, or if at
    least one motor command stays above the maximum of 255 for more than \SI{200}{ms}
    (saturation).
\end{description}
When balancing stops, all motors are switched off and braked, the run time and the cause are
reported, and the state returns to WAITING.

\section{Sensor processing}
\label{sec:sensor_processing}

\subsection{Edge angle}

On the edge, only the rotation about the sensor's X axis matters. The firmware uses the X tilt
angle $\varphi_x$ computed by the MPU6050\_light library, subtracts the calibrated offset
$\varphi_0$ and combines it with the gyroscope rate $\omega_x$ in a complementary filter:
\begin{equation}
  \theta_k = 0.98\left(\theta_{k-1} + \omega_x\,\Delta t\right) + 0.02\left(\varphi_x - \varphi_0\right).
  \label{eq:edge_cf}
\end{equation}
Here $\Delta t$ is the time since the last control step, not since the last filter step. The edge
gains were tuned with this behaviour, so it was deliberately left unchanged.

\subsection{Vertex tilt vector}

On the vertex, the tilt is a three-dimensional vector. The current gravity direction $\vect{n}$ is
the normalised accelerometer vector, $\vect{r}$ the normalised calibrated vertex pose. The
accelerometer-based tilt error (in degrees) is
\begin{equation}
  \vect{e}_\text{acc} = \frac{180}{\pi}\,\left(\vect{n}\times\vect{r}\right).
\end{equation}
The order of the cross product matters: $\vect{n}\times\vect{r}$ has the same sign as the
integrated gyroscope rate (with the MPU6050\_light conventions, a positive rotation about X
increases the measured Y acceleration). The opposite order was the main bug during development
(\cref{sec:history}).

The tilt estimate $\vect{e}$ is obtained with a complementary filter with the time constant
$\tau = \SI{0.5}{s}$:
\begin{equation}
  \vect{e}_k = a\left(\vect{e}_{k-1} + \vect{\omega}\,\Delta t\right) + (1-a)\,\vect{e}_\text{acc},
  \qquad a = \frac{\tau}{\tau + \Delta t},
  \label{eq:vertex_cf}
\end{equation}
where $\vect{\omega}$ is the gyroscope rate vector and $\Delta t$ the measured time since the last
filter step. The long time constant suppresses the disturbances of the accelerometer caused by the
wheel accelerations. After every filter step, the component of $\vect{e}$ along $\vect{r}$ is
removed, because a rotation about the vertical axis is not a tilt and is never corrected by the
accelerometer.

\section{Edge controller}
\label{sec:edge_controller}

The edge controller implements \cref{eq:edge_law} for motor X:
\begin{equation}
  u_X = \rho(t)\,\Bigl(k_{e1}\,(\theta - \theta_\text{trim}) + k_{e2}\,\omega_x + k_{e3}\,n_X\Bigr),
  \label{eq:edge_controller}
\end{equation}
with the gains $k_{e1}$, $k_{e2}$, $k_{e3}$ (\cmd{ek1}, \cmd{ek2}, \cmd{ek3}), the wheel speed
$n_X$ of motor X and the start-up ramp $\rho(t)$ that rises linearly from 0 to 1. The command
$u_X$ is limited to $\pm 255$.

\begin{lstlisting}[caption={Edge controller in \cmd{cubilox\_mk1.ino}.}, label={lst:edge}]
float output = (eK1 * edgeAngle + eK2 * gyroRateX + eK3 * motorSpeed_X) * rampFactor;
motorControl_X((int)output);
\end{lstlisting}

\subsection{Automatic trim}
\label{sec:trim}

If the calibrated balance point is slightly off, the wheel keeps turning in one direction
(\cref{sec:wheel_speed}). According to \cref{eq:equilibrium}, a positive wheel speed means that the
angle setpoint is too large. The trim therefore moves the setpoint slowly against the wheel speed:
\begin{equation}
  \theta_\text{trim} \leftarrow \theta_\text{trim} - k_\text{trim}\, n_X\, \Delta t ,
\end{equation}
limited to \SI{\pm 3}{\degree}, with the rate $k_\text{trim}$ (\cmd{trim}, default 0.001). The trim
is only active after the start-up ramp and is kept after a fall. On the real cube, the wheel
oscillates around zero and the trim stays within about \SI{\pm 0.1}{\degree}.

\section{Vertex controller}
\label{sec:vertex_controller}

The vertex controller works in the coordinate frame of the cube. The motors Y and Z are mounted
mirrored, which is described by the motor signs $s = (+1, -1, -1)$ and the encoder signs
$w = (+1, +1, +1)$ in \cmd{config.h}. The physical wheel momentum vector is
\begin{equation}
  \vect{h} = \bigl(s_X w_X n_X,\; s_Y w_Y n_Y,\; s_Z w_Z n_Z\bigr).
\end{equation}
Both the gyroscope rate $\vect{\omega}$ and the wheel momentum $\vect{h}$ are split into a tilt part
and a yaw part according to \cref{eq:split}. The torque command vector is
\begin{equation}
  \vect{t} = k_{v1}\,(\vect{e} - \vect{e}_\text{trim}) + k_{v2}\,\vect{\omega}_\text{tilt}
           + k_{v3}\,\vect{h}_\text{tilt}
           + \bigl(k_{yd}\,\omega_\text{yaw} - k_{yw}\,h_\text{yaw}\bigr)\,\vect{r},
  \label{eq:vertex_controller}
\end{equation}
where $\omega_\text{yaw} = \vect{\omega}\cdot\vect{r}$ and $h_\text{yaw} = \vect{h}\cdot\vect{r}$.
The gains are \cmd{vk1}, \cmd{vk2}, \cmd{vk3} for the tilt and \cmd{vyd}, \cmd{vyw} for the yaw
axis. The three motor commands are
\begin{equation}
  u_i = \rho(t)\, s_i\, t_i ,\qquad i \in \{X, Y, Z\}.
\end{equation}
The first three terms of \cref{eq:vertex_controller} are the three-dimensional version of the edge
controller. The yaw term damps the rotation about the vertical axis and slowly removes the yaw
momentum of the wheels; $k_{yw}$ must be positive, otherwise it is a positive feedback
(\cref{sec:yaw_physics}).

\begin{lstlisting}[caption={Core of the vertex controller in \cmd{cubilox\_mk1.ino}.},
  label={lst:vertex}]
float yawTorque = vYawD * gYaw - vYawW * hYaw;
float tX = vK1 * (cfVertexAngleX - vTrimX) + vK2 * gTx + vK3 * hTx + yawTorque * rx;
float tY = vK1 * (cfVertexAngleY - vTrimY) + vK2 * gTy + vK3 * hTy + yawTorque * ry;
float tZ = vK1 * (cfVertexAngleZ - vTrimZ) + vK2 * gTz + vK3 * hTz + yawTorque * rz;

float outX = MOTOR_SIGN_X * tX * rampFactor;
float outY = MOTOR_SIGN_Y * tY * rampFactor;
float outZ = MOTOR_SIGN_Z * tZ * rampFactor;
\end{lstlisting}

The vertex has an automatic trim as well, which works like the edge trim but as a vector
perpendicular to $\vect{r}$:
$\vect{e}_\text{trim} \leftarrow \vect{e}_\text{trim} - k_\text{vtrim}\,\vect{h}_\text{tilt}\,\Delta t$,
limited to \ang{3}. Its rate \cmd{vtrim} is 0 (off) by default, because the cube balances best
with a careful calibration by hand (\cref{sec:history}).

\section{Calibration and persistent storage}
\label{sec:calibration}

\subsection{Gyroscope}

The gyroscope has a small offset that changes with temperature. It is therefore measured at every
start: the cube only has to lie still for about three seconds. The surface does not need to be
level, because only the gyroscope is calibrated.

\subsection{Edge and vertex poses}

The balance points of the edge and the vertex are measured once with the commands \cmd{edge} and
\cmd{vertex} and stored in flash. The user balances the cube by hand on the edge or vertex;
after \SI{7}{s} (one short tick per second), the acceleration vector is averaged for \SI{3}{s}.
The measurement is rejected if
\begin{itemize}
  \item the standard deviation of one acceleration component exceeds \SI{0.05}{g} (the cube was
        moved), or
  \item the measured direction is more than \ang{12} away from an ideal pose: for the vertex a
        space diagonal $(\pm 1, \pm 1, \pm 1)/\sqrt{3}$, for the edge a diagonal in the Y--Z plane,
        i.e.\ an edge parallel to the axis of motor X.
\end{itemize}

Balancing by hand finds the \emph{real} balance point, where the centre of mass is exactly above
the pivot. A fixture that holds the cube geometrically straight does not, because the centre of
mass is generally not exactly in the geometric centre. This difference was the key to stable vertex
balancing (\cref{sec:history}).

The accelerometer offsets are not recomputed at start-up. The offsets that were active during the
calibration are stored together with the poses and loaded again at every start, so that the
stored poses remain valid after every restart and power loss. (A constant accelerometer offset does
not matter for balancing, because the poses are measured with the same offset.)

\begin{table}[htbp]
  \centering
  \caption{Data stored in the flash memory (ESP32 \cmd{Preferences}).}
  \label{tab:flash}
  \begin{tabularx}{\textwidth}{llX}
    \toprule
    Namespace & Content & Deleted by \\
    \midrule
    \cmd{cube}    & controller gains (after \cmd{save}) & \cmd{reset} \\
    \cmd{cuberef} & edge and vertex poses, accelerometer offsets & new calibration (overwrites) \\
    \cmd{cubebat} & battery voltage factor (after \cmd{vbatcal}) & new \cmd{vbatcal} \\
    \bottomrule
  \end{tabularx}
\end{table}

\section{Communication and logging}

Commands can be sent over the USB serial port or over Bluetooth Classic (serial port profile) with
any serial terminal; they are listed in \cref{sec:commands}. All output goes to both interfaces.

During development, an analysis of the Bluetooth library showed that it sends every print call
as a separate packet; with a connected client, its transmit queue holds 32 packets and a full
queue can block the caller for up to one second. Since the diagnostic output consisted of many
print calls, this could stall the control loop. The firmware therefore formats each log line completely and
sends it as one packet (\cmd{logLine()}), and it always writes the motor commands \emph{before}
any log output.

In normal operation, the cube only reports the start and the end of a balancing run (with the run
time and the cause) and a short status line every \SI{10}{s}. With \cmd{debug 1}, detailed lines
are printed (\cref{tab:debug}).

\begin{table}[htbp]
  \centering
  \caption{Debug output (\cmd{debug 1}).}
  \label{tab:debug}
  \begin{tabularx}{\textwidth}{lX}
    \toprule
    Line & Fields \\
    \midrule
    \cmd{IDLE} (every \SI{500}{ms}) & detected pose, acceleration vector, edge angle, vertex tilt,
      battery voltage \\
    \cmd{EDGE} (every \SI{500}{ms}) & time \cmd{t}, \cmd{angle}, wheel speed \cmd{wheel}, \cmd{trim},
      longest control gap \cmd{loop} \\
    \cmd{RUN} (every \SI{50}{ms}) & time \cmd{t}, tilt vector \cmd{e}, gyroscope \cmd{g}, wheel
      speeds \cmd{w}, motor commands \cmd{o}, yaw and tilt momentum, trim vector, \cmd{loop} \\
    \cmd{STOP} & tilt, acceleration magnitude and number of short shocks at the end of a vertex run \\
    \bottomrule
  \end{tabularx}
\end{table}

\section{Where to change what}

Most settings are constants in \cmd{config.h}:
\begin{itemize}
  \item \textbf{Pins} (\cmd{PWM\_X}, \cmd{DIR\_X}, \ldots): adapt them if the board is wired
        differently.
  \item \textbf{Default gains} (\cmd{eK1}--\cmd{eK3}, \cmd{vK1}--\cmd{vK3}, \cmd{vYawD},
        \cmd{vYawW}, trim rates): they are used until gains are stored with \cmd{save}.
  \item \textbf{Motor and encoder signs} (\cmd{MOTOR\_SIGN\_*}, \cmd{WHEEL\_SIGN\_*}): must match
        the mounting direction of the motors; the command \cmd{motortest} recommends them.
  \item \textbf{Start and stop limits} (\cmd{EDGE\_START\_ANGLE\_DEG},
        \cmd{VERTEX\_STOP\_ANGLE\_DEG}, ramp times, \cmd{SATURATION\_STOP\_MS}, \ldots).
  \item \textbf{Calibration timing} (\cmd{CALIB\_SETTLE\_MS} = \SI{7}{s},
        \cmd{CALIB\_AVERAGE\_MS} = \SI{3}{s}) and its plausibility limits.
  \item \textbf{Battery} (\cmd{VBAT\_FACTOR}, \cmd{BAT\_WARN\_V}, \cmd{BAT\_MIN\_V}).
  \item \textbf{Bluetooth name} (\cmd{BT\_DEVICE\_NAME}).
\end{itemize}
